\section{Research questions}\label{sec:rq}

\subsection{From the task to research questions}
The project brief asks for a memory over one long video that supports adding,
updating, merging, compressing, deleting and searching records; for questions
about past events, changes to objects, and relations between different parts
of the video; and for an analysis of what the system keeps and loses, how it
updates memory, how it finds relevant records, how accurate its answers are,
and whether its timestamps match the source clips. A comparison against a plain
text summary or another memory organisation, with the same observations and a
similar budget, is optional. \Cref{tab:task-map} shows where each requirement
is addressed.

\begin{table}[ht]
\centering\small
\caption{Requirements of the brief and where they are addressed.}\label{tab:task-map}
\begin{tabularx}{\linewidth}{@{}XX@{}}
\toprule
Requirement & Addressed by \\
\midrule
Memory structure; add, update, merge, compress, delete, search &
  structured memory C3 (\cref{sec:stage4}) \\
Questions on past events, object changes, cross-part relations &
  question types \emph{past event}, \emph{object change}, \emph{cross segment}
  (\cref{sec:stage2}) \\
What the system keeps and loses & RQ2 \\
How it updates memory after new events & RQ3 \\
How it finds relevant records & retrieval recall (RQ4) \\
Answer accuracy; timestamps vs.\ source clips & RQ1, RQ4 \\
Optional comparison, same observations and budget & baselines C1, C2, C4 (RQ1) \\
\bottomrule
\end{tabularx}
\end{table}

\subsection{Research questions}

\begin{description}[style=nextline,leftmargin=2.2em]
\item[RQ1 (accuracy and organisation).]
How accurately does the structured memory (C3) answer each type of question,
and how does it compare, at an equal budget $B$, with a rolling summary (C1),
a flat log (C2) and the unbounded log (C4)?
\item[RQ2 (what is kept and lost).]
Which information survives under a budget and which is lost first, and how
does accuracy change with evidence age $\age(q)$?
\item[RQ3 (updating).]
How reliably does the memory recognise that two mentions denote the same
entity and that an entity's attribute has changed, and how do errors in these
operations reach the answers?
\item[RQ4 (retrieval and grounding).]
Are the records relevant to a question retrieved, and do the records cited in
an answer have time support $I(C)$ overlapping the ground-truth evidence
$G(q)$?
\end{description}

\Cref{tab:rq-ops} lists how each question is measured. The analysis is
descriptive: every quantity is reported with a 95\% interval, and comparisons
between conditions are paired over questions.

\begin{table}[ht]
\centering\small
\caption{Operationalisation of the research questions.}\label{tab:rq-ops}
\begin{tabularx}{\linewidth}{@{}lXXl@{}}
\toprule
RQ & Primary measure & Secondary measures & Conditions \\
\midrule
RQ1 & $\Acc_m(B)$, overall and per question type & paired per-question
  differences (McNemar) & C1--C4, $B_1$--$B_3$ \\
RQ2 & evidence coverage (C2--C3); accuracy vs.\ $\age(q)$ (C1--C3) &
  operation-log trace of lost facts; memory size over time & C1--C3 \\
RQ3 & operation counts; merge precision and missed-merge rate (audit) &
  state-change precision; share of errors caused by corruption; ablations & C3 \\
RQ4 & retrieval recall; $g_m(q)$, $p_m(q)$, $\mathrm{tIoU}_m(q)$ & invalid citations & C2--C4 \\
\bottomrule
\end{tabularx}
\end{table}

\subsection{Expectations}
The design rests on four expectations, stated here so the reader can compare
them with the results; they are not tested as formal hypotheses.
\begin{itemize}[nosep]
\item Structured memory should help most on object-change and cross-segment
  questions, because state changes and entity identities are kept as explicit
  records that compression does not touch.
\item A rolling summary should lose fine attributes (colours, counts,
  intermediate states) before coarse events.
\item Differences between organisations should shrink as the budget grows.
\item Entity-resolution errors should be a major cause of failures on
  cross-segment questions.
\end{itemize}

\subsection{Failure attribution}\label{sec:attribution}
For each incorrect answer of a record-based organisation (C2--C4), exactly one
cause is assigned, checking the pipeline in order and stopping at the first
stage where the evidence is lost:
\begin{enumerate}[nosep]
  \item \emph{perception miss}: no observation $o_i$ with $c_i\cap G(q)\neq\varnothing$
        contains the evidence;
  \item \emph{storage loss}: it was observed but deleted or compressed away
        (traceable in the operation log);
  \item \emph{corruption}: it is stored, but a wrong merge (entity resolution)
        or a spurious state change distorted it;
  \item \emph{retrieval miss}: it is present in $M_n$ but not in $R(q,M_n)$;
  \item \emph{reasoning error}: it was retrieved, but the answer is wrong.
\end{enumerate}
The ordering makes the attribution unique and puts the burden of proof on the
earliest stage, which prevents perception errors from being misread as memory
failures.

\subsection{Statistical power and its consequences}\label{sec:power}
The planned evaluation set has 60--80 test questions. For an accuracy near
0.5 on $|\mathcal Q|=60$ questions, the half-width of a 95\% normal interval is
$1.96\sqrt{0.25/60}\approx 0.13$; for a single type with about 12 questions it
is $\approx 0.28$. Paired designs narrow these intervals considerably, because
all organisations answer the same questions, but differences below roughly
10 percentage points overall, and below 25 points within a type, cannot be
resolved. Two consequences follow. First, all comparisons are paired
(McNemar's test~\citep{mcnemar1947} for two organisations; percentile
bootstrap~\citep{efron1993bootstrap} over questions for intervals). Second,
every reported number carries a 95\% interval, and differences whose intervals
overlap substantially are described as not resolved rather than as effects.

\subsection{Threats to validity}\label{sec:threats}
\begin{itemize}[nosep]
  \item \emph{External validity}: one video, one genre. Findings concern this
        video; generalisation is argued, not measured.
  \item \emph{Perception ceiling}: all organisations inherit the errors of
        $\phi$. Its error rates are measured (Stage~3) and reported as the
        bound they impose.
  \item \emph{Annotator effects}: a single annotator writes and labels the
        questions. Mitigated by writing the questions before any memory
        exists and by tuning on the development split only.
  \item \emph{Judge bias}: free-form answers are scored with a rubric; any
        model-based judging is validated on a manually scored sample.
  \item \emph{Temporal resolution}: grounding is measured at chunk
        granularity $L$; claims about sub-chunk localisation are not possible.
\end{itemize}
